This class defines a program image (the executable associated with a
process). The only way to get a handle to a
\code{BPatch\_image} is via the \code{BPatch\_process} member
function \code{getImage}.

\begin{apient}
  bool findPoints(Dyninst::Address addr, std::vector<BPatch_point *> &points)
\end{apient}

\apidesc{
  Returns a vector of \code{BPatch\_point}s that correspond with the
  provided address, one per function that includes an
  instruction at that address. There will be one element if there is
  not overlapping code.
}

\begin{apient}
  std::vector<BPatch_variableExpr *> *getGlobalVariables()
\end{apient}

\apidesc{
  Return a vector of global variables that are defined in this image.
}

\begin{apient}
  BPatch_process *getProcess()
\end{apient}

\apidesc{
  Returns the \code{BPatch\_process} associated with this image.
}

\begin{apient}
  char *getProgramFileName(char *name, unsigned int len)
\end{apient}

\apidesc{
  Fills provided buffer name with the program's file name up to len
  characters. The filename may
  include path information.
}

\begin{apient}
  bool getSourceObj(std::vector<BPatch_sourceObj *> &sources)
\end{apient}

\apidesc{
  Fill sources with the source objects that belong to this image. If
  there are no
  source objects, the function returns false. Otherwise, it returns true.
}

\begin{apient}
  std::vector<BPatch_function *> *getProcedures(bool incUninstrumentable = false)
\end{apient}

\apidesc{
  Return a vector of the functions in the image. If the
  incUninstrumentable flag is set, the returned table of procedures
  will include uninstrumentable functions. The default behavior is to
  omit these functions.
}

\begin{apient}
  void getObjects(std::vector<BPatch_object *> &objs)
\end{apient}

\apidesc{
  Fill in a vector of objects in the image.
}

\begin{apient}
  std::vector<BPatch_module *> *getModules()
\end{apient}

\apidesc{
  Return a vector of the modules in the image.
}

\begin{apient}
  bool getVariables(std::vector<BPatch_variableExpr *> &vars)
\end{apient}

\apidesc{
  Fills vars with the global variables defined in this image. If
  there are no variable, the function returns false.
  Otherwise, it returns true.
}

\begin{apient}
  std::vector<BPatch_function*> *
  findFunction(const char *name, std::vector<BPatch_function*> &funcs,
               bool showError = true, bool regex_case_sensitive = true,
               bool incUninstrumentable = false)
\end{apient}

\apidesc{
  Return a vector of \code{BPatch\_function}s corresponding to name,
  or NULL if the function does not exist. If name contains a
  POSIX-extended regular expression, and \code{dont\_use\_regex} is
  false, a regular expression search will be performed on
  function names and matching \code{BPatch\_function}s returned. If
  showError is true, then Dyninst will report an error via
  the \code{BPatch::registerErrorCallback} if no function is found.
  If the incUninstrumentable flag is set, the returned table of
  procedures will include uninstrumentable functions. The
  default behavior is to omit these functions.
  
  [NOTE: If name is not found to match any demangled function names
    in the module, the search is repeated as if name is a
    mangled function name. If this second search succeeds, functions
    with mangled names matching name are returned
  instead.]
}

\begin{apient}

  bool (*BPatchFunctionNameSieve) (
    const char *name,
    void* sieve_data
  )  
  
  std::vector<BPatch_function*>*
  findFunction(std::vector<BPatch_function*> &funcs, BPatchFunctionNameSieve bpsieve,
               void *sieve_data = NULL, int showError = 0,
               bool incUninstrumentable = false)
\end{apient}

\apidesc{
  Return a vector of \code{BPatch\_function} s according to the generalized
  user\-specified filter function bpsieve.
  
  This permits users to easily build sets of functions according to
  their own specific criteria.
  Internally, for each \code{BPatch\_function} \code{f} in the image,
  this method makes a call to \code{bpsieve(f.getName(), sieve\_data)}.
  The user-specified function \code{bpsieve} is responsible for taking
  the name argument and determining if it belongs in the output vector,
  possibly by using extra user-provided information stored in
  \code{sieve\_data}. If the name argument matches the desired
  criteria, bpsieve should return true. If it does not, bpsieve
  should return false.

  If the incUninstrumentable flag is set, the returned table of
  procedures will include uninstrumentable functions. The
  default behavior is to omit these functions.
}

\begin{apient}
  bool findFunction(Dyninst::Address addr, std::vector<BPatch_function *> &funcs)
\end{apient}

\apidesc{
  Find all functions that have code at the given address, addr.
  Dyninst supports functions that share code, so this
  method may return more than one \code{BPatch\_function}. These
  functions are returned via the funcs output parameter. This
  function returns true if it finds any functions, false otherwise.
}

\begin{apient}
  BPatch_variableExpr *findVariable(const char *name, bool showError = true)
  BPatch_variableExpr *findVariable(BPatch_point &scope, const char *name)
\end{apient}

\apidesc{
  second form of this method is not implemented on Windows.
  Performs a lookup and returns a handle to the named variable. The
  first form of the function looks up only variables
  of global scope, and the second form uses the passed
  \code{BPatch\_point} as the scope of the variable. The returned
  \code{BPatch\_variableExpr} can be used to create references (uses)
  of the variable in subsequent snippets. The scoping rules
  used will be those of the source language. If the image was not
  compiled with debugging symbols, this function will
  fail even if the variable is defined in the passed scope.
}

\begin{apient}
  BPatch_type *findType(const char *name)
\end{apient}

\apidesc{
  Performs a lookup and returns a handle to the named type. The
  handle can be used as an argument to
  \code{BPatch\_addressSpace::malloc} to create new variables of the
  corresponding type.
}

\begin{apient}
  BPatch_module *findModule(const char *name, bool substring_match = false)
\end{apient}

\apidesc{
  Returns a module named name if present in the image. If the match
  fails, NULL is returned. If \code{substring\_match} is
  true, the first module that has name as a substring of its name is
  returned (e.g. to find libpthread.so.1, search for
  libpthread with \code{substring\_match} set to true).
}

\begin{apient}
  bool getSourceLines(unsigned long addr, std::vector<BPatch_statement> & lines)
\end{apient}

\apidesc{
  Given an address addr, this function returns a vector of pairs of
  filenames and line numbers at that address. This
  function is an alias for \code{BPatch\_­process::getSourceLines}.
}

\begin{apient}
  bool getAddressRanges(const char *fileName, unsigned int lineNo,
                        std::vector<std::pair< unsigned long, unsigned long>> &ranges )
\end{apient}

\apidesc{
  Given a file name and line number, fileName and lineNo, this
  function returns a list of address ranges that this source
  line was compiled into. This function is an alias for
  \code{BPatch\_process::getAddressRanges}.
}

\begin{apient}
  bool parseNewFunctions(std::vector<BPatch_module*> &newModules,
                         const std::vector<Dyninst::Address> &funcEntryAddrs)
\end{apient}

\apidesc{
  This function takes as input a list of function entry points
  indicated by the funcEntryAddrs vector, which are used to
  seed parsing in whatever modules they are found. All affected
  modules are placed in the newModules vector, which
  includes any existing modules in which new functions are found, as
  well as modules corresponding to new regions of the
  binary, for which new \code{BPatch\_modules} are created. The
  return value is true in the event that at least one previously
  unknown function was identified, or false otherwise.
}
